% Equations, figures and tables are numbered per appendix section: (A.1),
% Figure A.1, Table A.1.
\appendix
\numberwithin{equation}{section}
\counterwithin{figure}{section}
\counterwithin{table}{section}
%
\section{Test elements and meshes}
\label{app:elements}
%
The single test elements are defined by 4- or 8-node input elements (corner nodes counter-clockwise, followed by the mid-side nodes of the edges 1--2, 2--3, 3--4 and 4--1) and lifted to degree $p$ by isoparametric interpolation of the input geometry:
\begin{center}
  \small
  \begin{tabular}{@{}ll@{}}
    \toprule
    element & nodes $(x,y)$ \\
    \midrule
    parallelogram & $(0,0)$, $(2,0.3)$, $(2.6,1.5)$, $(0.6,1.2)$ \\
    trapezoid & $(0,0)$, $(2,0)$, $(1.4,1)$, $(0.2,1.3)$ \\
    curved & $(0,0)$, $(2,0)$, $(2,1)$, $(0,1)$; $(1,-0.2)$, $(2.25,0.5)$, $(1,1)$, $(0,0.5)$ \\
    general & $(0,0)$, $(2,0.1)$, $(1.5,1.2)$, $(0.2,1)$; $(1.1,-0.15)$, $(1.9,0.6)$, $(0.8,1.25)$, $(0.05,0.5)$ \\
    \bottomrule
  \end{tabular}
\end{center}
The meshes of \cref{sec:num-convergence} consist of $N\times N$ 8-node elements on the unit square whose interior mid-side nodes are moved normal to their edge by $\delta h$ with $h=1/N$, so that the relative sagitta of every interior edge is $\delta$. The normal is oriented consistently, so that all interior horizontal edges bow in the $-y$-direction and all interior vertical edges in the $-x$-direction. The $3\times2$ meshes of \cref{sec:num-patch,sec:num-spectra} on $[0,3]\times[0,2]$ use the same curved-edge distortion with $\delta=0.15$ and $0.20$ (with $h=1$). For the mesh denoted by \emph{angular + curved $0.30$}, the interior corner nodes are first shifted alternately by $\pm0.3$ in the $x$-direction, before the curved-edge distortion with $\delta=0.15$ is applied.
%
\section{Derivations}
\label{app:derivations}
%
For the sake of completeness, this appendix collects the derivations of the statements used in \cref{sec:formulation,sec:basis,sec:hybrid}. All functions of an element are considered as functions on the reference square $\Kh$, composed with the geometry map $\bx_K$. We write $\Qp=\Qp_p(\xi,\eta)$, $\Pp_p(\xi,\eta)$ for the polynomials of total degree $p$ in the reference coordinates, $\hat\Pp_p=\{q\circ\bx_K:q\in\Pp_p(x,y)\}$ for the complete polynomials of degree $p$ in $(x,y)$ and $\hat\Qp=\{q\circ\bx_K:q\in\Qsk\}$ for the trial space of \PM. $\Ihat$ denotes the GLL interpolation onto $\Qp$, and $\mathcal I_{\mathrm{T}}$ and $\mathcal I_{\mathrm{Q}}$ denote the nodal interpolation by the trial functions of \PH\ and \PM. For anisotropic degrees, $p$ is replaced by $m=\min(p_\xi,p_\eta)$ in $\Pp_p$ and $\hat\Pp_p$, and all arguments carry over.

\paragraph{Skew coordinates under affine transformations (\cref{sec:skew}).} If the nodes are transformed by $\bx\mapsto\matr A\bx+\vect c$ with $\det(\matr A)\ne0$, the geometry map becomes $\bx_K'=\matr A\bx_K+\vect c$. Hence, $\matr J_0'^\top=\matr A\matr J_0^\top$ and $\bx'-\bx_0'=\matr A(\bx-\bx_0)$, so that $\bxib'=\matr J_0^{-\top}\matr A^{-1}\matr A(\bx-\bx_0)=\bxib$. The bounding box \cref{eq:bbox} is computed from the skew coordinates of the nodes and is therefore unchanged as well.

\paragraph{Basis invariance (\cref{sec:basis-invariance}).} For another basis $\bm\psi=\matr G\bm\phi$ with a regular matrix $\matr G$, we obtain $\matr V_\psi=\matr G\matr V_\phi$ and $\matr P_\psi=\matr G\matr P_\phi$, so that $\matr V_\psi^{-1}\matr P_\psi=\matr V_\phi^{-1}\matr G^{-1}\matr G\matr P_\phi=\matr V_\phi^{-1}\matr P_\phi$. The functions $M_i^K=\sum_c(\matr V^{-1})_{ic}\phi_c$ lie in the trial space and satisfy $M_i^K(\bx_j)=\delta_{ij}$. If $\matr V$ is regular, they are the only such functions. For the nodal data $\vect f=\matr V^\top\vect c$ of a function of the trial space with the coefficient vector $\vect c$, the first-order error of the computed interpolant is $\vect c^\top\delta\matr V\,\matr N_{\mathrm M}$, where $\delta\matr V$ denotes the backward error of the factorization. Up to a constant, this error is bounded by $u\,\|\vect c\|\,\|\matr V\|\,\Lam_{\mathcal U}$. For polynomials of low degree, $\|\vect c\|\,\|\matr V\|$ is moderate in every basis, whereas generic nodal data can have $\|\vect c\|\approx\|\matr V^{-1}\|\,\|\vect f\|$.

\paragraph{Condition number in the Lagrange basis, \cref{eq:kappa}.} With the Lagrange basis, $\phi_c(\bm g_d)=\delta_{cd}$. Choosing the query points $\bm g_d$ in \cref{eq:vandermonde} gives $\matr P=\matr I$ and hence $\matr N_{\mathrm M}=\matr V^{-1}$, which gives $(\matr V^{-1})_{id}=M_i^K(\bm g_d)$. The 1-norm is the largest absolute column sum. Thus, $\|\matr V\|_1=\max_j\sum_c|\phi_c(\bxit_j)|=\max_j\lambda_{\mathrm{ref}}(\bxit_j)$ and $\|\matr V^{-1}\|_1=\max_d\sum_i|M_i^K(\bm g_d)|=\max_d\lambda_{\mathcal U}(\bm g_d)$.

\paragraph{Orthogonal complement (\cref{sec:hybrid-definition}).} Let $a+b>p$ and $i+j\le p$. Then $i<a$ or $j<b$, and the orthogonality of $L_a$ to all polynomials of degree less than $a$ gives $\int_{\Kh}L_a(\xi)L_b(\eta)\,\xi^i\eta^j\dd\xi\dd\eta=\int L_a\xi^i\dd\xi\int L_b\eta^j\dd\eta=0$. The functions $L_aL_b$ with $a,b\le p$ are linearly independent and span $\Qp$, those with $a+b\le p$ span $\Pp_p(\xi,\eta)$, and $\dim\Cp_p=n-\dim\Pp_p$. Hence, $\Cp_p$ is the $L^2(\Kh)$-orthogonal complement of $\Pp_p(\xi,\eta)$ in $\Qp$, which is defined independently of any basis.

\paragraph{Properties of the hybrid trial space (\cref{sec:hybrid-properties}).} (i) The hybrid space $\hat\Tp_p=\hat\Pp_p+\Cp_p$ is spanned by $n$ functions whose values at the $n$ nodes form the regular matrix $\matr V_{\mathrm{T}}$. Hence, they are linearly independent, the sum is direct, and the nodal map is bijective. (ii) and (iii) follow immediately from the definition and the invariance of the skew coordinates and of $\bxi$. (iv) For a bilinear map, $x,y\in\Qp_1(\xi,\eta)$, and hence $x^ay^b\circ\bx_K\in\Qp_{a+b}(\xi,\eta)\subset\Qp$ for $a+b\le p$. With $\Cp_p\subset\Qp$, we obtain $\hat\Tp_p\subset\Qp$, and equality follows from (i).

\paragraph{Representation of the interpolants, \cref{eq:repT}.} Let $\pi$ denote the projector of $\Qp$ onto $\Pp_p(\xi,\eta)$ along $\Cp_p$, i.e., the truncation of the Legendre expansion to $a+b\le p$, and define $S=\pi\Ihat|_{\hat\Pp_p}:\hat\Pp_p\to\Pp_p(\xi,\eta)$ and $S_{\mathrm{Q}}=\Ihat|_{\hat\Qp}:\hat\Qp\to\Qp$. We assume that $\bx_K$ has a point with $\det(\matr J)\ne0$, so that $\dim\hat\Pp_p=\dim\Pp_p$. The nodal values of a function $\hat v$ are in one-to-one correspondence with $\Ihat\hat v\in\Qp$, and $\Ihat c=c$ for $c\in\Cp_p\subset\Qp$. Let $\hat v=q+c$ with $q\in\hat\Pp_p$ and $c\in\Cp_p$ have the nodal values of $f$, i.e., $\Ihat q+c=\Ihat f$. Applying $\pi$, whose kernel is $\Cp_p$, gives $Sq=\pi\Ihat f$. Conversely, for every $q$ with $Sq=\pi\Ihat f$, the function $c=\Ihat f-\Ihat q$ lies in $\ker\pi=\Cp_p$. If $S$ is invertible, $q$ and $c$ are therefore uniquely determined for every $f$, and $\hat v=q+\Ihat f-\Ihat q$, which is \cref{eq:repT} with $q=S^{-1}\pi\Ihat f$. If $S$ is singular, we take $0\ne q\in\ker S$ and $c=-\Ihat q\in\Cp_p$. The function $\hat v=q+c$ then vanishes at all nodes. Either $\hat v\ne0$, or $q=-c\in\hat\Pp_p\cap\Cp_p$, so that the sum is not direct and $\dim\hat\Tp_p<n$. In both cases, the hybrid shape functions do not exist. Thus, the hybrid space is unisolvent at the nodes if and only if $S$ is invertible. Likewise, an interpolant $\hat v\in\hat\Qp$ satisfies $\Ihat\hat v=\Ihat f$, i.e., $S_{\mathrm{Q}}\hat v=\Ihat f$, and $\hat v=\Ihat\hat v+(I-\Ihat)\hat v$, so that $\mathcal I_Qf=\Ihat f+(I-\Ihat)S_{\mathrm{Q}}^{-1}\Ihat f$. With $\|\Ihat f\|_\infty\le\Lam_{\mathrm{N}}\|f\|_\infty$, we obtain the bounds
%
\begin{equation}
  \Lam_{\mathrm{T}}\le\Lam_{\mathrm{N}}\bigl(1+\|(I-\Ihat)S^{-1}\pi\|\bigr)\qquad\text{and}\qquad
  \Lam_{\mathrm{Q}}\le\Lam_{\mathrm{N}}\bigl(1+\|(I-\Ihat)S_{\mathrm{Q}}^{-1}\|\bigr),
  \label{eq:lebbound}
\end{equation}
%
with operator norms on $\Qp$ in the maximum norm on $\Kh$. For the nodal data of node $i$, $\Ihat f=\hat N_i$, so that the shape functions are $\hat M_{i,\mathrm{T}}=\hat N_i+(I-\Ihat)S^{-1}\pi\hat N_i$ and $\hat M_{i,\mathrm{Q}}=\hat N_i+(I-\Ihat)S_{\mathrm{Q}}^{-1}\hat N_i$. If $\hat\Pp_p\subset\Qp$, in particular for every bilinear map, $(I-\Ihat)q=0$, $\mathcal I_{\mathrm{T}}=\Ihat$ and $\Lam_{\mathrm{T}}=\Lam_{\mathrm{N}}$.

\paragraph{Factorization of the hybrid Vandermonde matrix, \cref{eq:factor}.} Let the rows of $\matr V_{\mathrm{ref}}$ be the nodal values of $L_a(\xi)L_b(\eta)$, $a,b\le p$, ordered into the $n_p$ low-order rows $\matr L$ ($a+b\le p$) and the rows $\matr R$ of the complement. The hybrid Vandermonde matrix is ordered alike, $\matr V_{\mathrm{T}}=[\matr Q;\matr R]$, where the rows of $\matr Q$ are the nodal values of a basis $\{q_k\}$ of $\hat\Pp_p$. Every row of $\matr Q$ has a unique expansion in the rows of the regular matrix $\matr V_{\mathrm{ref}}$, whose coefficients are the Legendre coefficients of $\Ihat q_k$, i.e., $\matr Q=\matr S^\top\matr L+\matr B\matr R$, where $\matr S$ is the matrix of $S$ and $\matr B$ contains the coefficients of the complement. Hence, $\matr V_{\mathrm{T}}=\bigl[\begin{smallmatrix}\matr S^\top&\matr B\\\matr0&\matr I\end{smallmatrix}\bigr]\bigl[\begin{smallmatrix}\matr L\\\matr R\end{smallmatrix}\bigr]$, and \cref{eq:factor} follows. On an affine element, $q\circ\bx_K\in\Pp_p(\xi,\eta)$, $\Ihat$ is exact, and $\matr S$ is the matrix of the invertible composition with an affine map. Since $\det(\matr S)$ depends continuously on the nodes, $\matr S$ is regular on every element that is sufficiently close to an affine one.

\paragraph{Interpolation estimates, \cref{eq:interp}.} Both trial spaces reproduce $\Pp_p(x,y)$, so that $f-\mathcal I_{\mathcal U}f=(f-v)-\mathcal I_{\mathcal U}(f-v)$ for every $v\in\Pp_p(x,y)$. We choose $v$ as the Taylor polynomial of degree $p$ of $f$ about a point of the convex hull $K^\ast$ of $K$. Then, $\max_{K^\ast}|f-v|\le Ch_K^{p+1}|f|_{p+1}$ and $\max_{K^\ast}|\nabla(f-v)|\le Ch_K^{p}|f|_{p+1}$, with $C$ depending only on $p$. The first estimate in \cref{eq:interp} follows from $\max_K|\mathcal I_{\mathcal U}(f-v)|\le\Lam_{\mathcal U}\max_K|f-v|$ and the second from $|\nabla\mathcal I_{\mathcal U}(f-v)|\le\max_K|f-v|\sum_i|\nabla M_i^K|\le h_K^{-1}\Lam'_{\mathcal U}\max_K|f-v|$. Under a similarity transformation $\bx\mapsto\alpha\matr R\bx+\vect c$, the shape functions are unchanged as functions of the material point, their gradients scale with $\alpha^{-1}$ and $h_K$ with $\alpha$, so that $\Lam_{\mathcal U}$ and $\Lam'_{\mathcal U}$ are invariant.
